\section{Slow proportions with changing domains}
\label{sec:proportions}

We continue the construction at one fixed star $v$, using the parameters
of Lemmas~\ref{lem:parameter-gaps} and~\ref{lem:caps}.
The caps $d_i$ will be reduced to fit the parent amount.  This section
constructs proportions for those reductions.  The construction depends
only on the star's observed rank posteriors and held parameters; the
parent amount $q_v$ is not an input.

Write $I$ for the active children, and extend all vectors by zero to the
fixed set of children when comparing different times.  Recall that
$B_i$ is the sum of the current rank probabilities over the core flags,
$S=\sum_i B_i$, and inactive children have no core flags.
If a dominant child $o$ is marked, the other active children are called
sides.  For regular children and for a marked child with sides,
respectively, the held parameters satisfy
\begin{equation}\label{eq:proportion-parameter-facts}
\begin{gathered}
  1\le h_i\le C\ell\quad(i\text{ regular}),\qquad
  h_o=u=U/A\in(0,1),\qquad H=1+\log(1/u),\\
  h_i\ge cH\quad(i\text{ a side}),\qquad
  \sum_{i\text{ regular}}e^{-h_i}\le
  \begin{cases} C,&\text{no dominant child},\\ Cu,&\text{a dominant child}.
  \end{cases}
\end{gathered}
\end{equation}
The last assertion follows from $e^{-h_i}\le a_i/A$ and the comparison
of $U$ with the side total.  These facts hold throughout each epoch.
When $I$ is empty there is no proportion to choose.  When $I=\{o\}$,
we use $\alpha_o=1$, $u=h_o=0$, and no side vector.  Both potentials
below are then defined to be zero.  In particular, we never evaluate a
denominator with $u=0$.

\subsection{Removal constraints and the simplex tracker}

Fix a sufficiently small absolute $c_t>0$ and put
$\eta_i=c_t h_i/\ell$.  For a nonempty active domain we require
\begin{equation}\label{eq:alpha-feasibility}
  \alpha_i\ge0,\qquad \sum_{i\in I}\alpha_i=1,
  \qquad S\alpha_i\le(1+\eta_i)B_i\quad(i\in I).
\end{equation}
This is automatic on a singleton domain.

To see the role of the tolerance $\eta_i$, temporarily fix any parent
amount $q_v\ge L_v^{\beta_v}$.  Removing the excess
$R_0=(\sum_i d_i-q_v)_+$ in proportions $\alpha$ would give
$q_i=d_i-R_0\alpha_i$.  The cap bounds and superadditivity imply
\[
 R_0\le\beta_v S+e_o,\qquad
 0\le e_o:=b_o-f_o^{\beta_v}\le\epsilon_b f_o^{\beta_v}/\ell,
\]
where $e_o=0$ without a marked child.  For $S>0$, a regular child's
removal is therefore at most
$(\beta_v+e_o/S)(1+\eta_i)B_i$.  The term $\beta_v B_i$ is
already allowed by the core decomposition
$L_i^{\beta_v}=n_i-\beta_v B_i+f_i^{\beta_v}$.
The parameter gap $\beta_i-\beta_v$ pays the extra coefficient,
provided $S$ is at least a fixed multiple of the dominant flex mass,
which is comparable to $D$.  At the marked child the error $e_o$
itself can be subtracted from its upper cap, since $\alpha_o\le1$.

When $S$ is small relative to $D$, the ratio $e_o/S$ need not be
controlled.  We then choose side removals $\omega_i$ first and give
the marked child the remaining parent amount, capped by $d_o$.
Superadditivity bounds its uncapped remaining amount below by
\[
             L_o^{\beta_v}
                  +\sum_{i\ne o}(\omega_i-\beta_v B_i).
\]
It therefore suffices to bound the total signed shortfall of the side
removals below $\beta_v B_i$ by a small multiple of $uD/\ell$; the
marked child's flex slack pays this error.  The removals must also stay below
$\beta_iB_i$ to preserve the side lower bounds.  The side tracker
below constructs $w$ so that $\omega_i=Dw_i$ has both properties.
Section~\ref{sec:allocation-finish} verifies the two output rules,
including their zero cases, and chooses between them by overlapping
comparisons of $S$ and $D$.

Near dominance also changes the simplex tracker itself.  The marked
child's parameter gap can be only of order $u/\ell$, forcing
$\eta_o=c_tu/\ell$.  Its potential must remain effective as
$u\downarrow0$; the different denominator below ensures this.
We will also normalize its primitive at $\alpha_o=1$, so changing
$u$ costs in proportion to the side sizes rather than the full
dominant size.

On every nonsingleton domain set, for $y\ge0$,
\begin{equation}\label{eq:proportion-denominators}
  F_i(y)=\frac{(h_i+y)^2}{h_i}\quad(i\text{ regular}),
  \qquad F_o(y)=u+Hy(1+y)\quad(o\text{ marked}).
\end{equation}

\begin{lemma}[Uniform denominator estimates]\label{lem:denominators}
For every nonsingleton domain above,
\begin{equation}\label{eq:denominator-estimates}
  \int_0^\infty\frac{dy}{F_i(y)}\le C,
  \qquad
  \sum_{i\in I}\alpha_i F_i\bigl(\log(1/\alpha_i)\bigr)
       \le C\sum_{i\in I}\alpha_i h_i
\end{equation}
for every $\alpha$ in its unit simplex.  A summand at $\alpha_i=0$
means its continuous extension, which is zero.  All constants are
uniform as $u$ tends to zero.
\end{lemma}

\begin{proof}
The regular integral equals $1$.  For the marked child, split at $1$:
\[
 \int_0^1\frac{dy}{u+Hy(1+y)}
   \le \frac{\log(1+H/u)}H,
 \qquad
 \int_1^\infty\frac{dy}{u+Hy(1+y)}\le\frac1H.
\]
Since $\log(1/u)=H-1$ and
$\log(1+H/u)\le H-1+\log(1+H)\le2H$, these bounds are absolute.

For the second assertion put $y_i=\log(1/\alpha_i)$ at positive
coordinates.  If $i$ is regular and $y_i\le4h_i$, its contribution
is at most $25\alpha_i h_i$.  If $y_i>4h_i$, then $h_i\ge1$ and
\[
 e^{-y_i}\frac{(h_i+y_i)^2}{h_i}
       \le C e^{-y_i}y_i^2
       \le C e^{-y_i/2}\le C e^{-h_i}.
\]
Thus the regular tail is at most $C$ without a marked child and at
most $Cu$ with one.  Without a marked child,
$\sum_i\alpha_i h_i\ge1$, which pays this tail.

With a marked child use the elementary inequality
$e^{-y}y(1+y)\le C(1-e^{-y})$, valid for all $y\ge0$.
It follows, including the endpoint $\alpha_o=0$, that
\[
 \alpha_o F_o\bigl(\log(1/\alpha_o)\bigr)
       \le C\{\alpha_o u+H(1-\alpha_o)\}.
\]
By \eqref{eq:proportion-parameter-facts},
$\sum_i\alpha_i h_i$ bounds this last expression up to an absolute
factor.  It also bounds $u$ up to an absolute factor: its dominant
contribution is $\alpha_o u$, and its side contribution is at least
$cH(1-\alpha_o)\ge cu(1-\alpha_o)$.  This pays the regular tail
and completes the proof.
\end{proof}

For held parameters and a nonnegative vector $B$, with
$S=\sum_i B_i$, define
\begin{equation}\label{eq:alpha-primitive}
\begin{aligned}
 g_i^\alpha(a;B)
   &=C_\alpha\ell\int_0^a
       \frac{S-(1+\eta_i)B_i/s}{F_i(\log(1/s))}\,ds
       &&(i\text{ regular}),\\
 g_o^\alpha(a;B)
   &=C_\alpha\ell\int_1^a
       \frac{S-(1+\eta_o)B_o/s}{F_o(\log(1/s))}\,ds,
       &&(o\text{ marked}),\\
 G^\alpha(\alpha;B)&=\sum_{i\in I}g_i^\alpha(\alpha_i;B).
\end{aligned}
\end{equation}
The improper integrals at zero converge by
$y=\log(1/s)$ and Lemma~\ref{lem:denominators}.
Consequently these functions extend continuously to the whole
simplex.  No convexity in $\alpha$ is asserted or needed.

We record the bounds needed when parameters change.  Uniformly in
$\alpha$,
\begin{equation}\label{eq:alpha-input-slopes}
 |G^\alpha(\alpha;B)-G^\alpha(\alpha;\widetilde B)|
       \le C\ell\|B-\widetilde B\|_1.
\end{equation}
Here and below changing $B$ also changes $S$ to its sum.  Indeed the
coefficient of $B_i$ in its own $B_i/s$ integral is bounded by
$C\ell\int_0^\infty dy/F_i(y)$.  The sum of the coefficients of
$S$ in the regular components has absolute value at most
$C\ell\sum_i\alpha_i\le C\ell$, because $F_i\ge h_i\ge1$.
The marked component's $S$ coefficient is bounded by
$C\ell\int_{\alpha_o}^1 ds/(sF_o(\log(1/s)))\le C\ell$.
This proves \eqref{eq:alpha-input-slopes} without a factor depending
on the number of active children.

The component bounds are
\begin{equation}\label{eq:alpha-component-bounds}
\begin{aligned}
 |g_i^\alpha(\alpha_i;B)|
   &\le C\ell(S\alpha_i+B_i) &&(i\text{ regular}),\\
 |g_o^\alpha(\alpha_o;B)|
   &\le C\ell\left(S(1-\alpha_o)+B_{\rm side}+uS\right),
   & B_{\rm side}&=\sum_{i\ne o}B_i.
\end{aligned}
\end{equation}
For the second bound, let $L=\log(1/\alpha_o)$, allowing $L=\infty$.
The change of variables in \eqref{eq:alpha-primitive} gives
\[
 g_o^\alpha(\alpha_o;B)
  =-C_\alpha\ell\int_0^L
       \frac{S(e^{-y}-1)+B_{\rm side}-\eta_o B_o}{F_o(y)}\,dy.
\]
For $0\le y\le L$, the absolute numerator is at most
$S(1-\alpha_o)+B_{\rm side}+\eta_o S$.
Since $\eta_o=c_tu/\ell$, Lemma~\ref{lem:denominators} proves
the bound.  The base point $1$ is responsible for this cancellation;
it permits parameter changes to be charged to side sizes even when
the dominant size is much larger.

\begin{lemma}[Simplex update]\label{lem:alpha-tracker}
Choose $C_\alpha$ sufficiently large after fixing $c_t$.
For any old vector $\alpha^0$ in the same simplex, a global minimizer
\begin{equation}\label{eq:alpha-minimization}
 \alpha^1\in\argmin_\alpha
       \left\{G^\alpha(\alpha;B)+S\|\alpha-\alpha^0\|_1\right\}
\end{equation}
satisfies \eqref{eq:alpha-feasibility} and
\begin{equation}\label{eq:alpha-update-payment}
 S\|\alpha^1-\alpha^0\|_1
       \le G^\alpha(\alpha^0;B)-G^\alpha(\alpha^1;B).
\end{equation}
The same assertions hold with the automatic singleton rule.
\end{lemma}

\begin{proof}
Continuity and compactness give a global minimizer, and comparison
with $\alpha^0$ gives \eqref{eq:alpha-update-payment}.
If $S=0$, then every $B_i=0$, the potential vanishes identically,
and feasibility is automatic.  We may retain $\alpha^0$.
Suppose henceforth that $S>0$ and the domain is nonsingleton.
For a positive coordinate write
\[
 \gamma_i=\frac{C_\alpha\ell}{F_i(\log(1/\alpha_i))}
                  \left(S-\frac{(1+\eta_i)B_i}{\alpha_i}\right)
\]
for the derivative of its primitive.
If $\alpha_i=0$ and $B_i>0$, the inward derivative is $-\infty$:
the negative term has size proportional to
$1/(sF_i(\log(1/s)))$, which tends to infinity as $s\downarrow0$.

There must be a coordinate with derivative less than $-2S$, with
this convention at zero.  Otherwise all zero coordinates have
$B_i=0$, and the inequalities $\gamma_i\ge-2S$ give
\begin{equation}\label{eq:alpha-derivative-contradiction}
 1\le\sum_{\alpha_i>0}\frac{\alpha_i}{1+\eta_i}
       \left(1+\frac{2F_i(\log(1/\alpha_i))}{C_\alpha\ell}\right).
\end{equation}
There is a fixed bound $\eta_i\le\eta_{\max}$, so, with
$T_\alpha=\sum_i\alpha_i h_i$,
\[
 \sum_i\frac{\alpha_i}{1+\eta_i}
    \le1-\frac{c_t}{1+\eta_{\max}}\frac{T_\alpha}{\ell}.
\]
The extra sum in \eqref{eq:alpha-derivative-contradiction} is at
most $CT_\alpha/(C_\alpha\ell)$ by
\eqref{eq:denominator-estimates}.  Choose $C_\alpha$ so this is
less than half the displayed deficit.  Since $T_\alpha>0$ on
every nonsingleton domain, this is a contradiction, however small
$u$ and $T_\alpha$ may be.

If a coordinate violates \eqref{eq:alpha-feasibility}, it is
positive and has strictly positive primitive derivative.  Transfer
a small amount $\varepsilon$ from it to a coordinate whose
derivative is below $-2S$.  These coordinates are distinct and the
transfer stays in the simplex.  The primitive sum decreases at a
rate greater than $2S$, whereas the movement penalty can increase
by at most $2S\varepsilon$.  At a zero receiving coordinate the
one-sided derivative $-\infty$ gives the same strict improvement
for sufficiently small $\varepsilon$.  This contradicts global
minimality and proves feasibility.
\end{proof}

\subsection{The side tracker}

When a dominant child has sides, choose a sufficiently small
absolute $c_w>0$ and define
\[
 \theta_i=\frac{c_w h_i}{\ell},\qquad
 z_i=\frac{a_i}{Ah_i}\quad(i\text{ a side}).
\]
Both numbers are positive.  Fix an absolute $C_w$ larger than
$2\sup_i(\beta^++\theta_i)$, with an additional positive margin.
The supremum is bounded absolutely by
\eqref{eq:proportion-parameter-facts} and the fixed parameter interval.
The side domain is $w_i\ge0$, $\sum_iw_i\le C_w$.
We require
\begin{equation}\label{eq:side-feasibility}
\begin{aligned}
 Dw_i&\le(\beta^++\theta_i)B_i&&\text{always},\\
 Dw_i&\ge\beta^+B_i-\theta_i z_iD&&\text{if }S\le2D.
\end{aligned}
\end{equation}
The total additive slack has the useful bound
\begin{equation}\label{eq:side-slack}
 \sum_{i\text{ side}}\theta_i z_i
      =\frac{c_w}{\ell}\sum_{i\text{ side}}\frac{a_i}{A}
      \le \frac{Cc_wu}{\ell}.
\end{equation}
With no sides, $w$ is the zero vector and all these statements use
empty sums.  When no dominant child is marked, set $D=0$ for the
sole purpose of writing the side movement term.

For fixed side parameters define
\begin{equation}\label{eq:side-primitive}
\begin{aligned}
 g_i^w(w_i;B_i,D)
  &=C'_w\int_0^{w_i}
       \frac{Ds-(\beta^++\theta_i/2)B_i+D\theta_i z_i/2}
            {\theta_i(s+z_i)}\,ds,\\
 G^w(w;B,D)&=\sum_{i\text{ side}}g_i^w(w_i;B_i,D).
\end{aligned}
\end{equation}
These potentials are continuous and affine in $(B,D)$.
For explicit bounds, set $L_i=\log(1+w_i/z_i)$.  Integration gives
\begin{equation}\label{eq:side-primitive-explicit}
 g_i^w=C'_w\left\{
   \frac{D}{\theta_i}(w_i-z_iL_i)
   +\frac{Dz_iL_i}{2}
   -\frac{(\beta^++\theta_i/2)B_iL_i}{\theta_i}
                    \right\}.
\end{equation}
Since $a_i/A<1$ for a side, $h_i=1+\log(A/a_i)$, and therefore
\[
 \log(1+C_w/z_i)
   =\log\left(1+C_wh_i e^{h_i-1}\right)\le Ch_i.
\]
Also $0\le z_iL_i\le w_i$ and
$0\le w_i-z_iL_i\le w_i$.  It follows that
\begin{equation}\label{eq:side-potential-bounds}
\begin{gathered}
 |g_i^w(w_i;B_i,D)|\le C\ell(Dw_i+B_i),\\
 |G^w(w;B,D)-G^w(w;\widetilde B,\widetilde D)|
   \le C\ell\left(\|B-\widetilde B\|_1+|D-\widetilde D|\right).
\end{gathered}
\end{equation}
In the second line the coefficient of $D$ is bounded by
$C\ell\sum_iw_i\le C\ell$, and the coefficient of each $B_i$
is bounded by $C\ell$.  Thus no factor counts the sides.

\begin{lemma}[Side update]\label{lem:side-tracker}
Choose the absolute multiplier $C'_w$ sufficiently large after
fixing the small constants.  For every old vector $w^0$ in the
side domain, a global minimizer
\begin{equation}\label{eq:side-minimization}
 w^1\in\argmin_w\left\{G^w(w;B,D)+D\|w-w^0\|_1\right\}
\end{equation}
satisfies \eqref{eq:side-feasibility}, and
\begin{equation}\label{eq:side-update-payment}
 D\|w^1-w^0\|_1\le G^w(w^0;B,D)-G^w(w^1;B,D).
\end{equation}
This holds for arbitrary nonnegative $B,D$ with $S=\sum_iB_i$,
without a requirement that they arise from ranks.
\end{lemma}

\begin{proof}
Compactness gives existence and comparison with $w^0$ gives the
payment.  Suppose first that $D>0$, and abbreviate
$b=\beta^+$, $\theta=\theta_i$, $z=z_i$, and $s=w_i$.
At an upper violation, $B_i<Ds/(b+\theta)$.  The numerator in
\eqref{eq:side-primitive} is consequently greater than
\[
 \frac{D\theta}{2}\left(\frac{s}{b+\theta}+z\right).
\]
Because $b+\theta$ has an absolute upper bound, a large $C'_w$
makes the primitive derivative greater than $2D$.
The violating coordinate is positive and can be decreased while
remaining in the side domain.  The penalty changes by at most
$D$ times this decrease, a contradiction to minimality.

Having proved upper feasibility, assume $S\le2D$.  Then
\[
 \sum_iw_i\le \max_i(b+\theta_i)\frac{\sum_{i\text{ side}}B_i}{D}
           \le2\max_i(b+\theta_i)<C_w.
\]
Thus every coordinate has room for a small increase.  At a lower
violation, $B_i>D(s+\theta z)/b$.  Substitution makes the same
numerator strictly less than
\[
 -\frac{D\theta}{2}\left(\frac{s}{b}+z\right);
\]
The derivative is therefore less than $-2D$ when $C'_w$ is large.
Increasing this coordinate strictly improves the objective, again
a contradiction.  This proves the lower inequality.

If $D=0$, the upper inequality is automatic.  The lower inequality
is required only when $S=0$, in which case every $B_i=0$ and it
too is automatic.  The potential need not vanish when $D=0<S$,
but minimization and its payment remain valid.  If both $D=S=0$,
the potential vanishes, and retaining the old side vector is allowed.
The empty side domain consists of its unique zero vector.
\end{proof}

\subsection{Adaptive parameters and domains}

The preceding updates will run at every step on which they are
defined, regardless of which output rule is subsequently used.
We now justify their total movement under all the changes from
Section~\ref{sec:allocation}.

We use explicitly the following coarse charges, whose expected
total is at most $CE_v+D_v$ by \eqref{eq:coarse-data-budget}:
the number of changed core flags; the sum of old and new $a_i$
at changed size coordinates; the old and new total active sizes
at wholesale refreshes; the old and new side totals at changes of
$U$; and $|\Delta D|$ off wholesale steps.
Here $D_v$ is finite and may depend on the fixed star and $k$,
but is independent of the input law and horizon.  Absolute constants
may absorb the fixed coarse tolerances.  The finely tracked caps
$b_i$ are not arguments of either proportion potential.

\begin{lemma}[Adaptive proportion movement]\label{lem:adaptive-proportions}
There is a causal update schedule satisfying
\eqref{eq:alpha-feasibility} and \eqref{eq:side-feasibility} after
every step, such that
\begin{equation}\label{eq:proportion-movement}
 \E\sum_{t\text{ not wholesale}}
   \left(S_t\|\alpha_t-\alpha_{t-1}\|_1
        +D_t\|w_t-w_{t-1}\|_1\right)
       \le C\ell E_v+D'_v.
\end{equation}
Vectors are extended by zero to the fixed child set.  The side
term is zero without a dominant child.  The simplex term includes
the automatic singleton state; an empty active domain contributes
zero.  The finite constant $D'_v$ is uniform in the law and horizon.
\end{lemma}

\begin{proof}
We give the schedule and its entire accounting.  All comparisons
below are made on a finite set of child and rank indices.

\paragraph{Filtering and drift precede every update decision.}
At the start of step $t$, hold the old proportion states, old
parameters, old $D$, and old core flags.  First evaluate the
potentials with the filtered old hidden counts, and then with the
new hidden counts under the new conditioning, still using those
old flags and parameters.  With these quantities held, each
potential is affine in the rank probabilities: for $G^w$, its
term linear in the held $D$ is simply constant on these two legs.
The signed expected filtering change is therefore zero.  By
\eqref{eq:alpha-input-slopes} and \eqref{eq:side-potential-bounds},
the absolute drift change is at most $C\ell$ times the sum of
absolute child-rank drifts.  Summed over all steps, its expectation
is at most $C\ell E_v$.  These statements concern all steps
before selecting any refresh or domain-change events.

Denote the input at the end of this evaluation by $\overline B$:
it uses the final rank probabilities and the old flags.  Thus
$\overline B_i\le n_i^{\rm old}$ even if filtering made the
current ranks very different.  We now compute the new flags,
sizes, and all other held parameters by the rules already given.

\paragraph{Wholesale steps.}
At a wholesale step reinitialize the proportion domains, choose
any vector in each new compact domain, and apply the appropriate
minimizing updates; use the automatic rules for empty or singleton
domains.  The absolute old potential, evaluated at
$\overline B$ and old $D$, is at most $C\ell$ times the old
total active size.  To see this for $G^\alpha$, sum
\eqref{eq:alpha-component-bounds} and use
$\sum_i\overline B_i\le\sum_i n_i^{\rm old}\le
C\sum_i a_i^{\rm old}$.  For $G^w$, use
\eqref{eq:side-potential-bounds}, $\|w\|_1\le C_w$, and
$D^{\rm old}\le C\sum_i a_i^{\rm old}$.
The new potential is bounded by the corresponding new total.
Thus the entire potential jump at this step is bounded by its
coarse wholesale charge times $C\ell$.  No movement at this step
is included in \eqref{eq:proportion-movement}.

\paragraph{Preparing a non-wholesale domain change.}
On a non-wholesale step $A$, $\beta^+$, and the presence and
identity of the dominant mark are fixed.  Empty/nonempty active
transitions are wholesale, and a marked dominant coordinate
cannot disappear within its epoch.  Let $I^-$ and $I^+$ be the
old and new active domains.  Define the synthetic input on the
fixed child set by
\begin{equation}\label{eq:synthetic-core-input}
 B_i^*=\begin{cases}B_i^{\rm new},&i\in I^-\cap I^+,\\
                     0,&\text{otherwise},\end{cases}
 \qquad S^*=\sum_iB_i^*.
\end{equation}
If $K_t$ is the number of changed core flags, then
\begin{equation}\label{eq:synthetic-core-errors}
 \|\overline B-B^*\|_1\le K_t,
 \qquad \|B^{\rm new}-B^*\|_1\le K_t,
 \qquad 0\le S^{\rm new}-S^*\le K_t.
\end{equation}
For common indices, the first discrepancy is just the sum over
changed flags, and each rank probability lies in $[0,1]$.
For an index leaving the active domain, all its new flags are
zero, so every old flag there is changed.  For an entering index,
all its old flags were zero, so every new flag there is changed.
These observations prove all three bounds.

Replace $\overline B$ by $B^*$ and old $D$ by $D^{\rm new}$,
holding the old domains and their other parameters.  The input
slope bounds charge this potential change to
$C\ell(K_t+|D^{\rm new}-D^{\rm old}|)$.
Run the minimizing updates on the old domains with inputs
$(B^*,S^*,D^{\rm new})$.  Denote the resulting vectors by
$\alpha^1,w^1$.  Lemmas~\ref{lem:alpha-tracker} and
\ref{lem:side-tracker} apply because their feasibility proofs use
only nonnegativity of these inputs and the held parameter bounds.
They pay
\begin{equation}\label{eq:prepared-movement}
 S^*\|\alpha^1-\alpha^{\rm old}\|_1,
 \qquad D^{\rm new}\|w^1-w^{\rm old}\|_1
\end{equation}
by decreases of the respective potentials.

Next delete departing coordinates and insert zero coordinates at
entering indices.  If $S^*>0$, a departing coordinate has
$B_i^*=0$, so \eqref{eq:alpha-feasibility} forces
$\alpha_i^1=0$.  Deletion and zero insertion then preserve the
simplex and cause no movement.  If $S^*=0$, every $B_i^*=0$,
the alpha potential is zero, and we may redistribute the simplex
mass to any vector in the new nonempty domain at zero weighted
cost.  If $D^{\rm new}>0$, the upper inequality in
\eqref{eq:side-feasibility} forces $w_i^1=0$ at each departing
side.  If $D^{\rm new}=0$, deletion of side coordinates has zero
weighted cost, and their primitive components vanish because
$B_i^*=D^{\rm new}=0$ there.  Inserting a coordinate at zero
always gives zero component potential.  Write $\alpha^2,w^2$
for the vectors on the new domains after this operation.

\paragraph{Changing the held parameters.}
Continue to use $B^*$ and $D^{\rm new}$.  At common regular
coordinates the prepared state satisfies
\begin{equation}\label{eq:prepared-coordinate-bounds}
 S^*\alpha_i^2\le CB_i^*,\qquad
 D^{\rm new}w_i^2\le CB_i^*\quad(i\text{ a side}).
\end{equation}
The assertions also hold at inserted zero coordinates, and at
$S^*=0$ or $D^{\rm new}=0$ with the evident zero left side.
The regular parameters $h_i,\eta_i,F_i$ and, on sides,
$\theta_i,z_i$ change only if $a_i$ changes.  Bound the absolute
old and new component values using
\eqref{eq:alpha-component-bounds},
\eqref{eq:side-potential-bounds}, and
\eqref{eq:prepared-coordinate-bounds}.  Their jump is at most
$C\ell B_i^*\le C\ell n_i^{\rm new}\le C\ell a_i^{\rm new}$,
which is paid by the old-plus-new size charge at that coordinate.
Departing and inserted components have already been seen to give
zero extra jump.

For a change in $u$, the prepared alpha vector satisfies
\[
 S^*(1-\alpha_o^2)\le C B_{\rm side}^*,\qquad
 B_{\rm side}^*\le C\sum_{i\text{ new side}}a_i^{\rm new},
 \qquad S^*\le C A.
\]
The first inequality follows by summing prepared side feasibility
when $S^*>0$ and is automatic when $S^*=0$; the last uses the
new core-size bound and the epoch comparison of total size with
$A$.  Applying the dominant component bound on both sides of
the parameter change shows that its jump is at most
\[
 C\ell\left(B_{\rm side}^*+
                  (u^{\rm old}+u^{\rm new})S^*\right)
 \le C\ell\left(U^{\rm old}+U^{\rm new}\right).
\]
This is a side-total charge, not a dominant-size charge.
It includes the appearance of the first side or disappearance of
the last side.  To spell out these endpoints, on a singleton
domain the potential is zero.  If $S^*>0$, preparation forces
all mass onto $o$ when the common side input is zero, and the
dominant primitive at $\alpha_o=1$ is zero.  If $S^*=0$, the
whole alpha potential is zero.  Thus assigning the singleton
potential to be zero is compatible with the estimates on both
sides of either transition.  The side vector on an empty side
domain is its unique zero vector.

\paragraph{The final update and its movement weight.}
With the new domains and parameters now installed, replace $B^*$
by $B^{\rm new}$, paying at most $C\ell K_t$ by
\eqref{eq:synthetic-core-errors} and the input slope bounds.
Run the two minimizing updates once more, now starting at
$\alpha^2,w^2$ and using the final inputs.  This pays
$S^{\rm new}\|\alpha^{\rm new}-\alpha^2\|_1$ and
$D^{\rm new}\|w^{\rm new}-w^2\|_1$.

The triangle inequality converts these payments to the differences
between the final and old states.  The alpha vectors at all these
substeps have norm one, hence any two are at distance at most two.
The domain-transition movement has weight zero at $S^*$.
Writing $\delta=S^{\rm new}-S^*\ge0$, we obtain
\begin{equation}\label{eq:adaptive-weight-comparison}
\begin{aligned}
 S^{\rm new}\|\alpha^{\rm new}-\alpha^{\rm old}\|_1
 &\le S^*\|\alpha^1-\alpha^{\rm old}\|_1
      +S^{\rm new}\|\alpha^{\rm new}-\alpha^2\|_1+4\delta,\\
 D^{\rm new}\|w^{\rm new}-w^{\rm old}\|_1
 &\le D^{\rm new}\|w^1-w^{\rm old}\|_1
      +D^{\rm new}\|w^{\rm new}-w^2\|_1.
\end{aligned}
\end{equation}
The first error is at most $4K_t$.  The second inequality needs
no change of weight, since both side updates used $D^{\rm new}$.
These statements also cover zero weights and singleton states.

\paragraph{Telescoping and uniformity.}
Let $\mathcal G_t=G_t^\alpha+G_t^w$ after the final update.
On non-wholesale steps the sum of the paid minimizing movements
is at most $\mathcal G_{t-1}-\mathcal G_t$, plus the signed
old-parameter filtering/drift change and the parameter/input
charges just listed.  On wholesale steps the same inequality
holds with no movement term and with the wholesale charge.
Summing over \emph{all} steps retains the zero expected filtering
sum; it does not condition that assertion on an adaptively
selected subset of steps.  Expected drift costs at most
$C\ell E_v$, and all remaining changes cost at most $C\ell$
times the coarse charges.  Equation~\eqref{eq:adaptive-weight-comparison}
therefore gives \eqref{eq:proportion-movement}, apart from the
initial and terminal potential magnitudes.

Those magnitudes have a uniform finite bound.  At a valid state,
$\sum_iB_i\le\sum_i n_i\le Ck$, $D\le Ck$,
$\|\alpha\|_1\le1$, and $\|w\|_1\le C_w$.
The component bounds give $|\mathcal G_t|\le C\ell k$, including
arbitrarily small positive $u$.  Only the fixed collection of
rank and parameter trackers contributes the other finite
initialization allowances.  Their states persist through all
epochs and inactive intervals; no new allowance is assigned when
a label or coordinate returns.  Thus $D'_v$ is independent of
the number of steps, histories, refreshes, and domain changes.
The schedule uses only current observed posteriors and prior
states.  Global minimizers can be selected by fixed tie-breaking
on the compact finite-dimensional domains, so it is causal.
\end{proof}
